\section{Coordination-Free Frontiers}
\label{sec:frontiers}

We ask how much a specification can promise before coordination is needed,
holding the inputs, input order, and output order fixed.  A specification $T$
is \emph{stronger} than $S$ when $T\subsetneq S$.  This compares sets of rows
and is unrelated to $\OutputOrd$, which compares outputs.

\subsection{Remembering what may remain true}

For $A,B\subseteq O$, the \emph{Hoare preorder} is
        \begin{equation}
  A\HoareOrd B
  \quad\Longleftrightarrow\quad
  \forall a\in A\;\exists b\in B:\;a\OutputOrd b.
  \label{eq:lifted-output-order}
        \end{equation}
Fix a table $X\subseteq I\times O$, with fibers $X(i)$ possibly empty.  A
\emph{powerset relaxation} of $X$ is a map $A:I\to\mathcal P(O)$ with
$X(i)\HoareOrd A(i)$ for every $i$; it satisfies \emph{Relational
Monotonicity} (Definition~\ref{def:relational-monotonicity}) when
$i\InputOrd k$ implies $A(i)\HoareOrd A(k)$.  The \emph{remembered output} is
        \begin{equation}
  \widehat X(i)=\bigcup_{j\InputOrd i}X(j).
  \label{eq:monotone-hull}
        \end{equation}

\begin{theorem}[Remembered-output frontier]
\label{thm:monotone-hull}
$\widehat X$ is a powerset relaxation of $X$ satisfying Relational
Monotonicity, and every such relaxation $A$ satisfies
        \begin{equation}
  \widehat X(i)\HoareOrd A(i)
  \qquad\text{for every }i.
  \label{eq:hull-coverage}
        \end{equation}
If every $\widehat X(i)$ is finite and $M(i)$ is its set of maximal elements,
then $M(i)\HoareOrd\widehat X(i)\HoareOrd M(i)$, $M$ satisfies Relational
Monotonicity, and $M(i)$ is the least subset $Y\subseteq\widehat X(i)$ with
$\widehat X(i)\HoareOrd Y$.
\end{theorem}

\begin{proof}
$X(i)\subseteq\widehat X(i)$, and $i\InputOrd k$ gives $\widehat
X(i)\subseteq\widehat X(k)$.  For Equation~\ref{eq:hull-coverage}, let
$o\in X(j)$ with $j\InputOrd i$.  Relaxation gives $a\in A(j)$ with
$o\OutputOrd a$, and Relational Monotonicity gives $b\in A(i)$ with
$a\OutputOrd b$.

In a finite poset every element lies below a maximal one, so $\widehat
X(i)\HoareOrd M(i)$; the converse holds since $M(i)\subseteq\widehat X(i)$.
Then $i\InputOrd k$ gives $M(i)\HoareOrd\widehat X(i)\HoareOrd\widehat
X(k)\HoareOrd M(k)$.  If $Y\subseteq\widehat X(i)$ and $\widehat X(i)\HoareOrd
Y$, each $m\in M(i)$ lies below some $y\in Y\subseteq\widehat X(i)$, and
maximality gives $y=m$; so $M(i)\subseteq Y$.  Minimality does not use
finiteness.
\end{proof}

When $X=S$ is a specification, every $\widehat X(i)$ is nonempty, and
$i\mapsto\widehat X(i)$ is a functional specification with outputs in
$(\mathcal P(O)\setminus\{\varnothing\})/{\equiv_H}$, where $A\equiv_H B$ means
$A\HoareOrd B\HoareOrd A$.  It is monotone, hence consistent by
Theorem~\ref{thm:functional-calm}.  The theorem says it is the least such
relaxation, and $M$ is its least representative drawn from remembered
outputs.

Table~\ref{tab:lookup-lift} gives an example: a data item moves from node $a$ to
node $b$ as input grows from $i_0$ to $i_1$.  Output $o_a$ ($o_b$) says that a
lookup from $a$ ($b$) succeeds, and $o_{a\leadsto b}$ says that $b$ holds the
item and $a$ forwards to it, so $o_a,o_b\OutputOrd o_{a\leadsto b}$.

\begin{table}[t]
  \centering
  \small
  \caption{Original, remembered, and pruned fibers for the lookup example.}
  \label{tab:lookup-lift}
  \begin{tabular}{@{}c l l l@{}}
    \toprule
    Input & $X(i)$ & $\widehat X(i)$ & $M(i)$ \\
    \midrule
    $i_0$ & $\{o_a\}$ & $\{o_a\}$ & $\{o_a\}$ \\
    $i_1$ & $\{o_b,o_{a\leadsto b}\}$
          & $\{o_a,o_b,o_{a\leadsto b}\}$
          & $\{o_{a\leadsto b}\}$ \\
    \bottomrule
  \end{tabular}
\end{table}

\subsection{Persistent forwarding}

B-link trees keep a search starting at a split node correct by following a
right link, so the parent can be repaired later~\cite{lehman1981efficient}.  We
abstract the invariant.  Inputs are finite legal restructuring histories $h$,
ordered by prefix.  For a key $k$ and search position $p$, $(k,p)\in R(h)$ when
the key-directed search from $p$ follows a finite path to a current occurrence
of $k$.  Each restructuring step must satisfy:
\begin{enumerate}
  \item search positions keep stable identities;
  \item every old key-directed edge is replaced by a key-directed path; and
  \item every old occurrence of $k$ has a path to a current occurrence of $k$.
\end{enumerate}
Deletion and reclamation can violate (2) and (3) and are excluded.

\begin{proposition}[Persistent-forwarding monotonicity]
\label{prop:persistent-forwarding}
Under (1)--(3), $h\InputOrd h'$ implies $R(h)\subseteq R(h')$.  The functional
specification $h\mapsto R(h)$, ordered by inclusion, is therefore monotone and,
by Theorem~\ref{thm:functional-calm}, consistent.
\end{proposition}

\begin{proof}
By induction it suffices to treat one step.  Replace each edge of a valid search
path by the path that (2) provides, then append the path that (3) provides from
the old occurrence to a current one.
\end{proof}

Let $X(g)\subseteq R(g)$ be the claims whose search is exact (uses no
forwarding) at $g$, and order claims by equality, so that $\HoareOrd$ is
inclusion.  By Proposition~\ref{prop:persistent-forwarding}, $R$ is a powerset
relaxation of $X$ satisfying Relational Monotonicity, so
Equation~\ref{eq:hull-coverage} gives
        \begin{equation}
  \widehat X(h)\subseteq R(h):
  \label{eq:forwarding-hull}
        \end{equation}
every claim that was ever exact remains valid, witnessed by a forwarding path.
Here $X$ has empty fibers, which is why Theorem~\ref{thm:monotone-hull} is
stated for tables.

\subsection{Preserving independent local choices}

Fix a finite participant set $P$ with $|P|\ge 2$.  Coordinates are indexed by
participants.  For each $p$ let $(W_p,\le_p)$ be a poset and, for each input
$i$, let $V_p(i)\subseteq W_p$ be nonempty: the \emph{views} $p$ may expose at
$i$.  A \emph{partial tuple} assigns to each $p$ either a value in $W_p$ or
$\bot$; partial tuples are ordered coordinate-wise, with $\bot$ below every
value.  $x$ is a \emph{restriction} of a complete tuple $y$ when each $x_p$ is
$\bot$ or $y_p$.
\begin{itemize}
  \item The \emph{product interface} $S$ has a row $(i,x)$ for every
        restriction $x$ of every $y\in\prod_pV_p(i)$.
  \item A \emph{normalized strengthening} $T$ selects, at each input, a set of
        tuples in $\prod_pV_p(i)$ and has as rows exactly their restrictions.
        It is \emph{locally conservative} when for every $i$, $p$, and $v\in
        V_p(i)$ some selected tuple $z$ at $i$ has $z_p=v$.
  \item The \emph{owner deployment} (Definition~\ref{def:deployment}) of a
        specification over partial tuples lets $p$ emit exactly its rows whose
        tuple is $\bot$ outside coordinate $p$.  Participant $p$ is the sole
        owner of coordinate $p$ and owns no other.
\end{itemize}

\begin{theorem}[Independent-product frontier]
\label{thm:product-frontier}
\begin{enumerate}
  \item If views continue (for $i\InputOrd j$ and $v\in V_p(i)$ there is $w\in
        V_p(j)$ with $v\le_p w$), the product interface is monotone across
        partitions (Definition~\ref{def:monotone-across-partitions}) under the
        owner deployment, so every separable run family passes the survivor
        test (Proposition~\ref{thm:complete-cap-deployment}).
  \item If each $V_p(i)$ is an antichain, no strict, normalized, locally
        conservative strengthening $T\subsetneq S$ is monotone across
        partitions under its owner deployment.
\end{enumerate}
\end{theorem}

\begin{proof}
(1) Take an authorized participant run family and a common \indesc $j$ of its
current inputs.  We construct $w\in\prod_pV_p(j)$ that \outrel every emitted
output; the construction does not use separability.  If $p$ has a run ending at
current input $i_p$, its remainder table has a row $(i_p,u)$, and by
Lemma~\ref{lem:run-residual} $u$ \outrel every output $p$ emitted.  As $u$ is a
restriction of some $y\in\prod_qV_q(i_p)$, continuation gives $w_p\in V_p(j)$
with $y_p\le_p w_p$.  Otherwise pick any $w_p\in V_p(j)$.  Each output $p$
emitted is $\bot$ outside coordinate $p$, and its value there is at most
$u_p\le_p y_p\le_p w_p$.

(2) Strictness gives a row $(i,x)\in S\setminus T$, with $x$ a restriction of
some $y\in\prod_pV_p(i)$ that $T$ does not select at $i$.  By local
conservativity, the partial tuple $e_p$ that is $y_p$ at $p$ and $\bot$
elsewhere is a row of $T$ at $i$, authorized for $p$.  Let every participant
run from the canonical state at $i$ and emit $e_p$.  Under any partition, two
participants in different components emit, so the family is separable; take
$j=i$.  A common \outnoun $u\in T(i)$ is a restriction of a selected $z$, and
$y_p\le_p u_p=z_p$ with both in $V_p(i)$, so $y=z$, a contradiction.
\end{proof}

The hypotheses matter.  The deployment is essential: once some $V_p(i)$ has
two incomparable views, the product interface is not monotone, since those two
rows have no common \outnoun.  Sole
ownership is needed: let two replicas each expose configuration 0 or 1, and let
$T$ require them to agree, keeping $(0,0)$ and $(1,1)$.  This $T$ is strict and
locally conservative, so by part (2) it is not monotone across partitions; if
instead one participant owns both coordinates, $T$ is monotone across
partitions.  Neither part says whether a participant can emit its views during
a partition; that is Deployment Availability
(Definition~\ref{def:deployment-availability}), a property of the
implementation.

\paragraph{Causal register.}
Let $V_p(i)$ contain the register views of process $p$ (a source write for each
of $p$'s reads) that some legal causal execution of the causally closed history
$i$ explains, with each read owned by its process.  An expansion preserves the
causal past of completed reads, so views continue and part (1) applies.  This
uses causal legality only: a write that completed before an old read in real
time can be added without entering that read's causal past.  For the same
reason linearizability can remove a single-process causal view, so it is not a
locally conservative strengthening of the causal product and part (2) does not
apply over all histories.  The post-reconciliation execution of
Appendix~\ref{app:deployment} exhibits the cross-process obstruction directly.
